\documentclass[11pt,a4paper]{article}
\usepackage[margin=25mm]{geometry}
\usepackage{fontspec}
\setmainfont{TeX Gyre Pagella}
\usepackage{amsmath,amssymb,mathtools}
\usepackage{graphicx}
\usepackage{xcolor}
\usepackage[unicode,colorlinks=true,linkcolor=blue,urlcolor=blue]{hyperref}
\usepackage{bookmark}
\usepackage{enumitem}
\setlist{itemsep=3pt,topsep=5pt}
\setlength{\emergencystretch}{3em}
\newcommand{\tightlist}{\setlength{\itemsep}{0pt}\setlength{\parskip}{0pt}}
\newcommand{\passthrough}[1]{#1}
\newcommand{\pandocbounded}[1]{#1}
\makeatletter
\def\maxwidth{\ifdim\Gin@nat@width>\linewidth\linewidth\else\Gin@nat@width\fi}
\def\maxheight{\ifdim\Gin@nat@height>0.75\textheight0.75\textheight\else\Gin@nat@height\fi}
\makeatother
\setkeys{Gin}{width=\maxwidth,height=\maxheight,keepaspectratio}
\hypersetup{pdfauthor={OpenAI GPT-6.1 Sol, Codex, Ultra effort}}
\begin{document}
\section{Hilbert symbols and local
conics}\label{hilbert-symbols-and-local-conics}

\emph{Written by OpenAI GPT-6.1 Sol in Codex, Ultra effort, October
2026. Self-checked by the writing AI; no independent review is claimed.
Public domain (CC0).}

Kummer theory labels abelian extensions by radicals, while reciprocity
labels their automorphisms by elements of the ground field. The Hilbert
symbol records the action of the latter on the former. This produces a
perfect pairing, a norm criterion, and explicit tests for whether a
quadratic conic has a local point.

Initially let \(K\) be a nonarchimedean local field containing
\(\mu_n\), where \(n\geq1\) is prime to \(\operatorname{char}K\). In
characteristic zero, this allows \(n\) divisible by the residue
characteristic. We use the Kummer theorem proved in
\href{../hilberts-theorem-90-and-kummer-theory.html}{Hilbert's theorem
90 and Kummer theory}, finite local reciprocity from
\href{../local-reciprocity-and-norm-groups.html}{Local reciprocity and
norm groups}, and its explicit arithmetic normalization from
\href{../explicit-local-reciprocity-and-existence.html}{Explicit local
reciprocity and the existence theorem}. The quadratic norm calculations
over \(\mathbf Q_2\) in
\href{../abelian-ramification-conductors-and-hasse-arf.html}{Abelian
ramification, conductors and Hasse--Arf}, Exercise 2, will give the wild
quadratic formula. Real quadratic symbols are treated in section 4.

\textbf{Prerequisite proof availability.} The named results below
identify specific programme lessons. The
\href{https://kokunoyumeto.github.io/open-math-courses/courses/NT-CFT/proof-dependencies.html\#lesson-11}{prerequisite
record} distinguishes published proofs, supplied owner texts awaiting
publication, and missing full proofs. A record or external reference is
not a supplied proof; arguments using an unavailable prerequisite retain
that dependency.

\subsection{1. Definition and norm
identities}\label{definition-and-norm-identities}

For \(a,b\in K^\times\), choose \(\lambda\) with \(\lambda^n=b\). The
field \(L_b=K(\lambda)\) is finite cyclic: all roots are
\(\zeta\lambda\), \(\zeta\in\mu_n\subset K\), and its Galois group
embeds in \(\mu_n\). Define \[
(a,b)_n=\frac{(a,L_b/K)(\lambda)}{\lambda}\in\mu_n.
\tag{1}
\] The \textbf{first} argument enters reciprocity; the \textbf{second}
is put under the root. Replacing \(\lambda\) by \(\zeta\lambda\) has no
effect because \(\zeta\in K\). This fixes the convention even when
\(n>2\).

\subsubsection{Proposition 11.1. Algebraic
properties}\label{proposition-11.1.-algebraic-properties}

The symbol is multiplicative in each argument, factors through
\(K^\times/K^{\times n}\) in each, and satisfies \[
\begin{gathered}
(a,b)_n=1\ \Longleftrightarrow\
 a\in N_{L_b/K}L_b^\times,\\
(a,-a)_n=1,\qquad (a,1-a)_n=1\quad(a\ne1),\\
(a,b)_n(b,a)_n=1.
\end{gathered}
\tag{2}
\] The elements in each displayed symbol are nonzero.

\textbf{Proof.} Multiplicativity in \(a\) is the homomorphism property
of the norm residue symbol. For the second argument, put roots of
\(b,c\) in their finite abelian compositum. Restriction compatibility
says its reciprocity automorphism acts on their product as the product
of its actions. That product is an \(n\)th root of \(bc\), proving
\((a,bc)_n=(a,b)_n(a,c)_n\). An \(n\)th power in the first argument
gives an \(n\)th power in \(\mu_n\), hence 1; multiplying the second
argument by an \(n\)th power only multiplies its root by an element of
\(K\). Thus the pairing factors through both quotients.

An automorphism of \(L_b\) fixes \(\lambda\) exactly when it is
identity. Finite reciprocity identifies the identity kernel with the
norm group, proving the first equivalence.

We prove the two special norm identities without assuming \([L_b:K]=n\).
The separable polynomial \(T^n-b\) defines a finite algebra \[
A_b=K[T]/(T^n-b).
\] Factoring it into relatively prime irreducible polynomials and
applying the polynomial Chinese remainder theorem expresses \(A_b\) as a
product of fields. Every factor is \(K\)-isomorphic to \(L_b\), because
its root is some \(\zeta\lambda\) and generates the same field. The
determinant norm of an element of this product is the product of its
field norms, and therefore lies in \(N_{L_b/K}L_b^\times\) when the
element is invertible.

Multiplication by \(T\) has eigenvalues the \(n\) roots of \(T^n-b\), so
\[
N_{A_b/K}(-T)=-b,\qquad
N_{A_b/K}(1-T)=1-b.
\tag{3}
\] The first follows by multiplying their negatives, and the second by
evaluating the monic polynomial at 1. These elements are invertible when
their displayed norms are nonzero. Taking \(b=-a\) proves that \(a\) is
a norm from \(L_{-a}\); taking \(b=1-a\) proves that \(a\) is a norm
from \(L_{1-a}\). These are the middle identities in (2).

Finally expand \(1=(ab,-ab)_n\). In its first factor use \(-ab=(-a)b\),
and in its second use \(-ab=(-b)a\). Multiplicativity yields \[
1=(a,-a)_n(a,b)_n(b,-b)_n(b,a)_n
  =(a,b)_n(b,a)_n.
\] This proves the last identity. \(\square\)

For \(n=2\), inversion of a sign changes nothing, so the pairing is
symmetric. It need not be alternating: \((a,a)_n=(a,-1)_n^{-1}\), and
\((-1,-1)_2\) over \(\mathbf Q_2\) will be \(-1\).

\subsection{2. Nondegeneracy}\label{nondegeneracy}

\subsubsection{Theorem 11.2. The perfect Hilbert
pairing}\label{theorem-11.2.-the-perfect-hilbert-pairing}

The group \(B=K^\times/K^{\times n}\) is finite, and the pairing induces
an isomorphism \[
B\longrightarrow\operatorname{Hom}(B,\mu_n),\qquad
b\longmapsto\bigl(a\mapsto(a,b)_n\bigr).
\tag{4}
\] In particular it is nondegenerate in both arguments.

\textbf{Proof.} First verify finiteness, including the wild case. Put
\(s=v_K(n)\), which is finite since \(n\ne0\) in \(K\). If
\(v_K(t)>2s\), the polynomial \[
P(X)=X^n-(1+t)
\] has \(v_K(P(1))>2s\) and \(v_K(P'(1))=s\). The Newton estimates
proved in the descent lemma of the preceding reciprocity lesson apply
starting at 1: an increment of valuation \(r>s\) gives a next increment
of valuation at least \(2r-s\), while the derivative keeps valuation
\(s\). The iterates converge in \(K\) to a unit whose \(n\)th power is
\(1+t\). Hence sufficiently deep units are \(n\)th powers.

The unit group modulo a deep unit subgroup is finite, and valuation
modulo \(n\) has only \(n\) classes. Thus \(B\) is finite. It has
exponent dividing \(n\).

If \(b\) is nontrivial in \(B\), the field \(L_b\) is nontrivial. Finite
reciprocity is onto its Galois group, so some \(a\) gives an
automorphism that moves its generating root. Therefore \((a,b)_n\ne1\).
This proves injectivity in (4).

A finite abelian group of exponent dividing \(n\) has exactly as many
\(\mu_n\)-valued characters as elements. The finite character lemma in
lesson 3, section 4, proves this directly, together with character
separation, for every finite abelian group killed by \(n\). Consequently
the injection (4) is onto. If \(a\) pairs trivially with every \(b\), it
is annihilated by every character of \(B\), and character separation
gives \(a=1\) in \(B\). This proves the other nondegeneracy. \(\square\)

Thus being a norm from every radical extension of exponent dividing
\(n\) is equivalent to being an \(n\)th power in \(K\).

\subsection{3. The tame formula and its
direction}\label{the-tame-formula-and-its-direction}

Now assume the residue characteristic \(p\) does not divide \(n\). Write
\(q=|\kappa_K|\), and \[
a=\pi^\alpha u,\qquad b=\pi^\beta v,\qquad u,v\in\mathcal O_K^\times.
\] Reduction maps \(\mu_n\) injectively to \(\kappa_K^\times\). Indeed,
if \(\zeta=1+x\) is an \(n\)th root with \(v_K(x)>0\), the term \(nx\)
in \((1+x)^n-1\) has strictly smaller valuation than every higher term,
so cannot cancel unless \(x=0\). Hence \(n\mid q-1\).

\subsubsection{Proposition 11.3. The tame Hilbert
symbol}\label{proposition-11.3.-the-tame-hilbert-symbol}

The symbol is the unique root in \(\mu_n\) whose reduction is \[
\overline{(a,b)_n}
 =\left((-1)^{\alpha\beta}
          \frac{\bar v^{\,\alpha}}{\bar u^{\,\beta}}\right)^{(q-1)/n}
 =\left(\overline{(-1)^{\alpha\beta}b^\alpha/a^\beta}\right)^{(q-1)/n}.
\tag{5}
\]

\textbf{Proof.} An \(n\)th root of a unit \(v\) lies in a finite
unramified extension. Choose such an extension whose residue field
contains a root of \(X^n-\bar v\); its derivative is a unit there.
Lifting that residue root and applying the Newton iteration with
derivative valuation zero constructs a root of \(X^n-v\) in the complete
unramified field. Existence of the field with that residue extension is
Theorem 2.1 of
\href{https://kokunoyumeto.github.io/open-math-courses/courses/NT-CFT/proof-dependencies.html\#NT-LOC-07}{Unramified
and totally ramified extensions}.

The symbol of a unit acts trivially on an unramified extension, so
\((u,v)_n=1\). The symbol of \(\pi\) there is arithmetic Frobenius. For
\(z^n=v\), its action has ratio with reduction \[
\overline{\varphi(z)/z}=\bar z^{\,q-1}
          =\bar v^{(q-1)/n}.
\] Thus \((\pi,v)_n\) has that reduction. Antisymmetry gives
\(\overline{(u,\pi)_n}=\bar u^{-(q-1)/n}\). Finally \((\pi,-\pi)_n=1\)
implies \[
\overline{(\pi,\pi)_n}=(-1)^{(q-1)/n};
\] the inverse of this sign equals itself. Expanding the pairing on
\(\pi^\alpha u,\pi^\beta v\) proves (5), also for negative integer
exponents. Injectivity of reduction on \(\mu_n\) identifies the symbol
itself. \(\square\)

In particular, under arithmetic normalization the numerator is
\(b^\alpha\). Swapping the two arguments inverts the symbol; this
distinction is visible for orders greater than two.

\subsection{4. Quadratic formulas}\label{quadratic-formulas}

Use the notation \((a,b)_p\) for the quadratic symbol over
\(\mathbf Q_p\). For odd \(p\), let \((u/p)\) be the sign
\(\bar u^{(p-1)/2}\) in the residue field, identified with \(\{\pm1\}\);
this is the Legendre symbol of a unit.

For odd \(2\)-adic units define, modulo 2, \[
\epsilon(u)=\frac{u-1}{2},\qquad
\omega(u)=\frac{u^2-1}{8}.
\tag{6}
\] These integral expressions depend only on the odd residue modulo 8.

\subsubsection{Theorem 11.4. Explicit quadratic
symbols}\label{theorem-11.4.-explicit-quadratic-symbols}

For \(a=p^\alpha u,\ b=p^\beta v\), with units \(u,v\), one has \[
(a,b)_p=(-1)^{\alpha\beta(p-1)/2}
           (u/p)^\beta(v/p)^\alpha\quad(p\text{ odd}),
\tag{7}
\] and \[
(a,b)_2=(-1)^{\epsilon(u)\epsilon(v)
                       +\alpha\omega(v)+\beta\omega(u)}.
\tag{8}
\] Over \(\mathbf R\), \[
(a,b)_\infty=-1
\quad\Longleftrightarrow\quad a<0\text{ and }b<0.
\tag{9}
\]

\textbf{Proof.} Formula (7) is (5) for \(n=2,q=p\); the inverse of each
Legendre sign equals the sign.

For (8), the square-class calculation in Exercise 2 of the ramification
lesson proves that \(-1,2,5\) form a basis of the eight square classes
of \(\mathbf Q_2\). Its independently computed norm groups give all the
following pairings: \[
\begin{array}{c|ccc}
(\ ,\ )_2&-1&2&5\\ \hline
-1&-1&1&1\\
2&1&1&-1\\
5&1&-1&1
\end{array}.
\tag{10}
\] Here \((-1,-1)_2=-1\) because \(-1\) is not in \(1+4\mathbf Z_2\),
the unit norm subgroup from \(\mathbf Q_2(i)\). The units \(-1\) and 5
are norms as appropriate for the other entries: \(-1\) is a unit norm
from \(\mathbf Q_2(\sqrt2)\), and every unit is a norm from the
unramified field \(\mathbf Q_2(\sqrt5)\). Also \(2=N(2+\sqrt2)\), giving
\((2,2)_2=1\). The element 2 has odd valuation and is not a norm from
that unramified quadratic field, giving \((2,5)_2=-1\). Symmetry
supplies the reflected entries, and \((5,5)_2=1\) is again an unramified
unit norm.

For every odd unit, \[
u\equiv(-1)^{\epsilon(u)}5^{\omega(u)}
\pmod{\mathbf Q_2^{\times2}}.
\tag{11}
\] Check its four residues: \(1,3,5,7\) have exponent pairs
\((0,0),(1,1),(0,1),(1,0)\). Their proposed representatives have exactly
these residues modulo 8, whose equality is the square criterion. Express
\(a,b\) with (11), reduce \(\alpha,\beta\) modulo 2, and expand by (10).
The only negative basis pairings are the \((-1,-1)\) entry and the two
\(2,5\) entries. Their exponents are exactly those in (8).

For the real case, the quadratic reciprocity map on \(\mathbf R^\times\)
takes positive elements to identity and negative elements to complex
conjugation. Its kernel is the norm image from \(\mathbf C^\times\):
\(N(x+iy)=x^2+y^2\) gives precisely the positive real numbers. If
\(b>0\), its square root is real, so the symbol is 1. If \(b<0\), its
root generates \(\mathbf C\), and a negative \(a\) acts by conjugation
and changes the root's sign. This proves (9). It also shows the real
pairing is nondegenerate on the two square classes. Over \(\mathbf C\)
all elements are squares and the quadratic symbol is 1. \(\square\)

\subsection{5. The conic criterion}\label{the-conic-criterion}

\subsubsection{Proposition 11.5. Local points on a diagonal
conic}\label{proposition-11.5.-local-points-on-a-diagonal-conic}

For \(K=\mathbf Q_p\) or \(\mathbf R\), and \(a,b\in K^\times\), \[
(a,b)_K=1
\quad\Longleftrightarrow\quad
z^2=ax^2+by^2
\text{ has a solution }(x,y,z)\ne(0,0,0)\text{ in }K^3.
\tag{12}
\]

\textbf{Proof.} If \(b\) is a square, the symbol is 1 and
\((0,1,\sqrt b)\) is a nonzero solution. Otherwise \(L=K(\sqrt b)\) is
quadratic. The norm criterion, including the real norm calculation, says
the symbol is 1 exactly when \(a=s^2-bt^2\) for some \(s,t\in K\). Such
an equality supplies the conic point \((1,t,s)\).

Conversely, if a nonzero solution has \(x=0\), then \(z^2=by^2\);
nonzeroness forces \(y\ne0\), contradicting that \(b\) is nonsquare.
Thus \(x\ne0\), and division by \(x^2\) gives \[
a=(z/x)^2-b(y/x)^2,
\] which is the norm of \(z/x+(y/x)\sqrt b\). Since \(a\ne0\), this is a
nonzero norm, proving the symbol is 1. \(\square\)

This is a test over one completion. A global assertion that a rational
conic has a rational point requires an additional local-to-global
theorem.

\subsection{6. Exercises and complete
solutions}\label{exercises-and-complete-solutions}

\subsubsection{Exercise 1 --- easy}\label{exercise-1-easy}

Compute \((2,3)_3\), \((-1,-1)_2\), and \((2,5)_5\).

\textbf{Solution.} For \((2,3)_3\), the first argument is a unit, the
second has valuation 1, so (7) gives the Legendre sign \((2/3)=-1\). For
\((-1,-1)_2\), both valuations are zero and both \(\epsilon(-1)\) equal
1 modulo 2; (8) gives \(-1\). For \((2,5)_5\), formula (7) gives
\((2/5)=-1\), because the nonzero squares modulo 5 are 1 and 4.

\subsubsection{Exercise 2 --- medium}\label{exercise-2-medium}

Prove the conic criterion, including solutions with a zero coordinate.

\textbf{Solution.} If \(b\) is square, the point \((0,1,\sqrt b)\)
proves solubility and the radical extension is trivial. If \(b\) is
nonsquare, a nonzero point cannot have \(x=0\): that would give either a
square root of \(b\), or all coordinates zero. Divide its equation by
\(x^2\) to express \(a\) as a quadratic norm, so the symbol is 1.
Conversely write a nonzero norm \(a=s^2-bt^2\) and take
\((x,y,z)=(1,t,s)\). This reasoning allows \(t=0\) or \(s=0\) and
applies over both the nonarchimedean fields and \(\mathbf R\).

\subsubsection{Exercise 3 --- medium}\label{exercise-3-medium}

Prove the Steinberg relation \((a,1-a)_n=1\) for \(a\ne0,1\).

\textbf{Solution.} Put \(b=1-a\) and form \(A_b=K[T]/(T^n-b)\).
Separability and \(\mu_n\subset K\) make each field factor isomorphic to
\(L_b=K(b^{1/n})\). The determinant norm of \(1-T\) is \[
\prod_{\lambda^n=b}(1-\lambda)=1-b=a.
\] It is nonzero, so \(1-T\) is invertible in the product. Its component
norms multiply to \(a\), and their product is a single norm from
\(L_b\), by norm multiplicativity. The symbol's norm criterion proves
the relation. Using the product algebra accounts for the case where the
radical field has degree smaller than \(n\).

\subsubsection{Exercise 4 --- hard}\label{exercise-4-hard}

Derive the complete quadratic formula (8) from square classes and norm
groups.

\textbf{Solution.} The odd-unit square criterion identifies the square
classes with \(2^\alpha(-1)^e5^w\), for \(\alpha,e,w\in\{0,1\}\). For an
odd unit \(u\), its four residues modulo 8 give \(e=\epsilon(u)\),
\(w=\omega(u)\), as in (11). The exact quadratic norm groups from the
ramification lesson give the symmetric table (10), with negative entries
only at \((-1,-1),(2,5),(5,2)\).

If \(a\) has exponent vector \((\epsilon(u),\alpha,\omega(u))\) in the
ordered basis \((-1,2,5)\), and \(b\) has vector
\((\epsilon(v),\beta,\omega(v))\), multiplicativity gives their sign
exponent \[
\epsilon(u)\epsilon(v)+\alpha\omega(v)+\omega(u)\beta.
\] This proves (8) for all classes. Squares in either argument do not
affect the symbol, so the calculation applies to arbitrary integer
valuations and all \(2\)-adic units.

\subsection{Editable edition}\label{editable-edition}

The reading edition provides the complete LaTeX source of this lesson,
the cumulative course LaTeX and the editable source ZIP. The archive
contains all twenty-four Markdown lessons, complete LaTeX bodies,
original diagrams, metadata and reproduction instructions.

\subsection{References}\label{references}

The norm characterization and finite perfect pairing lead to the tame,
real and dyadic quadratic formulas. The product-algebra proof handles
local conics, including the split case. All displayed symbol
normalizations are proved in the lesson.

\begin{itemize}
\tightlist
\item
  \href{https://www.jmilne.org/math/CourseNotes/CFT.pdf}{J. S. Milne,
  Class Field Theory, version 4.03}.
\end{itemize}

The \href{../free-proofs.html}{proof guide} gives the lesson sequence
and the exact prerequisite record. External references accompany the
written arguments.

\end{document}
